\documentclass[11pt]{article}
% Source-PDF: 数据结构 DATA STRUCTURE/后缀数组.pdf
\usepackage{oipl-vn}

\title{Mảng hậu tố}
\author{}
\date{}

\begin{document}
\maketitle

\origtitle{Suffix array}
\sourcepath{data-structure/suffix-array.pdf}

\section{Thuật toán nhân đôi}

Ý tưởng chính là trước hết sắp xếp từng ký tự đơn lẻ. Sau đó, dựa trên kết quả
sắp xếp từng ký tự, ta dịch vị trí để thu được kết quả sắp xếp của ký tự đơn
lẻ kế tiếp; sắp xếp lại sẽ thu được thứ tự của các xâu con gồm hai ký tự.
Tiếp tục dịch vị trí rồi trộn kết quả sắp xếp để thu được thứ tự của các xâu
con gồm bốn ký tự, v.v. Vì vậy, độ dài xâu con cần xác định thứ tự tăng theo
hàm mũ. Kết hợp với sắp xếp cơ số, mỗi lượt có độ phức tạp $O(n)$, nên tổng
độ phức tạp thời gian là $O(n\log n)$.

\begingroup
\small
\begin{verbatim}
int wa[maxn],wb[maxn],wv[maxn],ws[maxn];
int cmp(int *r,int a,int b,int l) {
    return r[a]==r[b]&&r[a+l]==r[b+l];
}
void da(int *r,int *sa,int n,int m) {
    int i,j,p,*x=wa,*y=wb,*t;
    for(i=0;i<m;i++) ws[i]=0;
    for(i=0;i<n;i++) ws[x[i]=r[i]]++;
    for(i=1;i<m;i++) ws[i]+=ws[i-1];
    for(i=n-1;i>=0;i--) sa[--ws[x[i]]]=i;
    for(j=1,p=1;p<n;j*=2,m=p)
    {
        for(p=0,i=n-j;i<n;i++) y[p++]=i;
        for(i=0;i<n;i++) if(sa[i]>=j) y[p++]=sa[i]-j;
        for(i=0;i<n;i++) wv[i]=x[y[i]];
        for(i=0;i<m;i++) ws[i]=0;
        for(i=0;i<n;i++) ws[wv[i]]++;
        for(i=1;i<m;i++) ws[i]+=ws[i-1];
        for(i=n-1;i>=0;i--) sa[--ws[wv[i]]]=y[i];
        for(t=x,x=y,y=t,p=1,x[sa[0]]=0,i=1;i<n;i++)
            x[sa[i]]=cmp(y,sa[i-1],sa[i],j)?p-1:p++;
    }
    return;
}
\end{verbatim}
\endgroup

\begin{enumerate}
  \item Bốn mảng \texttt{wa}, \texttt{wb}, \texttt{wv}, \texttt{ws} là các
  mảng phụ trợ. Trong quá trình chạy chương trình, chúng có các ý nghĩa khác
  nhau; các mục sau sẽ giải thích cụ thể.

  \item Hàm \texttt{cmp} ở dòng 2--4 dùng để so sánh xem hai khóa có giống
  nhau hay không. Khi ban đầu so sánh thứ tự từng ký tự, khóa là mã ASCII của
  ký tự; về sau, khi đồng thời so sánh thứ tự của $2^k$ ký tự với $k\ge 1$,
  khóa là thứ hạng của $2^{k-1}$ ký tự. Nếu hai khóa giống nhau, tức hai ký tự
  hoặc hai xâu con hoàn toàn giống nhau, hàm trả về $1$; nếu không, hàm trả về
  $0$. Cụ thể, khi so sánh các xâu độ dài $2^k$, giá trị \texttt{l} là
  $2^{k-1}$, nghĩa là nửa trước có cùng thứ hạng và nửa sau cũng có cùng thứ
  hạng.

  \item Trong khai báo hàm, \texttt{r} là xâu ban đầu, \texttt{sa} là mảng hậu
  tố, \texttt{n} là độ dài xâu ban đầu cộng thêm một. Lý do là phía sau xâu
  ban đầu được thêm một số nhỏ hơn mọi ký tự, chẳng hạn $0$, chủ yếu để tiện
  cài đặt hàm. \texttt{m} là tham số so sánh của sắp xếp cơ số; nếu xâu chỉ
  gồm chữ cái và chữ số thông thường theo bảng ASCII thì đặt \texttt{m = 128}
  là đủ.

  \item Dòng 7--10 thực hiện sắp xếp cơ số: tách xâu thành $n$ ký tự độ dài
  $1$ rồi sắp xếp, cuối cùng lưu kết quả vào mảng \texttt{sa}. Vòng lặp
  \texttt{for} ở dòng 10 bắt đầu từ \texttt{n - 1}; mục đích là khi gặp các ký
  tự giống nhau, ký tự có vị trí trước hơn trong xâu cũng có thứ tự trước hơn.

  \item Trong điều kiện vòng lặp ở dòng 11, \texttt{j} mỗi lần được nhân đôi,
  nghĩa là mỗi lượt dùng kết quả đã sắp xếp của lượt trước làm nền. Biến
  \texttt{p} có hai tác dụng: một là chỉ số tăng dần của mảng, hai là thứ hạng
  của xâu con đứng sau cùng sau khi sắp xếp. Nếu \texttt{p} bằng \texttt{n},
  mọi xâu con đã có thứ tự khác nhau và vòng lặp kết thúc. Câu lệnh
  \texttt{m = p} là một tối ưu nhỏ: trong lần sắp xếp dùng thứ hạng của lượt
  trước làm khóa, tham số sắp xếp lớn nhất chỉ cần là \texttt{p}.

  \item Dòng 13--14 thực hiện sắp xếp theo khóa thứ hai. Chẳng hạn nếu độ dài
  của ký tự gốc là $1$, khóa thứ hai chính là ký tự ở vị trí kế tiếp. Vì khóa
  thứ hai của $j$ ký tự cuối là rỗng, hoặc bằng $0$, chúng chắc chắn nằm ở đầu
  thứ tự; dòng 13 xử lý bước này. Do đã có thứ tự theo khóa thứ nhất, chỉ cần
  dịch trái mỗi giá trị trong mảng \texttt{sa} một vị trí là thu được vị trí
  bắt đầu của các xâu con đã sắp theo khóa thứ hai.

  \item Dòng 15 lưu thứ hạng của các vị trí bắt đầu, vốn đã nhận được theo thứ
  tự khóa thứ hai, vào mảng \texttt{wv}. Mảng này được dùng trong lần sắp xếp
  cơ số tiếp theo để xác định thứ tự giữa hai xâu con. Khi chỉ có một ký tự,
  giá trị được lưu không phải là thứ hạng mà là giá trị ứng với ký tự đó, nhưng
  vai trò là như nhau.

  \item Dòng 16--19 tiếp tục sắp xếp các xâu con, vốn đã có thứ tự theo khóa
  thứ hai, theo khóa thứ nhất. Chú ý rằng mảng \texttt{wv} lấy giá trị từ
  \texttt{x}, vì \texttt{y} lưu chỉ số của các xâu con đã có thứ tự theo khóa
  thứ hai. Vòng lặp cuối bắt đầu từ \texttt{n - 1}, nhờ đó với hai khóa thứ
  nhất giống nhau, vị trí đứng sau trong xâu vẫn đứng sau trong thứ tự cuối
  cùng; nói cách khác, thứ tự theo khóa thứ hai vẫn được giữ ổn định.

  \item Dòng 20--21 tìm thứ hạng tương ứng cho các xâu con đã có thứ tự, với
  thông tin thứ tự đang nằm trong mảng \texttt{sa}. Cần chú ý rằng hai xâu con
  hoàn toàn giống nhau phải nhận cùng một thứ hạng. Mảng \texttt{x} thu được
  cuối cùng lưu thông tin thứ hạng của từng xâu con.

  \item Sau một số lượt lặp, thứ tự của mọi xâu con đều khác nhau, vòng lặp
  kết thúc. Khi đó \texttt{sa} là mảng hậu tố, còn \texttt{x} là mảng thứ hạng.
\end{enumerate}

Tác giả gốc nói rằng phải tự mô phỏng theo ví dụ mới hiểu rõ thuật toán. Ví dụ
trong bài dùng xâu \texttt{aabaaaab}; khi cài đặt, ta xét thêm ký tự canh
\texttt{0} ở cuối:
\[
  r=\{97,97,98,97,97,97,97,98,0\},\qquad n=9.
\]
Sau lượt sắp xếp từng ký tự, ta có
\[
  sa=[8,0,1,3,4,5,6,2,7],
  \qquad
  x=[1,1,2,1,1,1,1,2,0].
\]
Ở lượt $j=1$, sau khi sắp xếp theo khóa thứ hai rồi khóa thứ nhất:
\[
  y=[8,7,0,2,3,4,5,1,6],
  \qquad
  sa=[8,0,3,4,5,1,6,7,2],
\]
và mảng thứ hạng trở thành
\[
  x=[1,2,4,1,1,1,2,3,0].
\]
Ở lượt $j=2$:
\[
  y=[7,8,6,1,2,3,4,5,0],
  \qquad
  sa=[8,3,4,5,0,6,1,7,2],
\]
cuối cùng
\[
  x=[4,6,8,1,2,3,5,7,0].
\]
Vì đã có \texttt{p = n}, mọi hậu tố đã được phân biệt và thuật toán kết thúc.

\section{Thuật toán DC3}

Ý tưởng chính là chia toàn bộ xâu thành hai phần: các hậu tố có vị trí bắt đầu
đồng dư $0$ theo modulo $3$, ký hiệu là $A$, và các hậu tố có vị trí bắt đầu
không đồng dư $0$ theo modulo $3$, ký hiệu là $B$. Sau đó nối hai phần của
$B$, tức phần đồng dư $1$ và phần đồng dư $2$, để tạo thành một xâu mới. Cứ
mỗi $3$ ký tự của xâu mới được xem như một ký tự rồi đem sắp xếp cơ số để thu
được thứ tự. Nếu vẫn tồn tại hai ký tự hoàn toàn giống nhau, ta gọi đệ quy hàm
này cho đến khi mọi ký tự đều phân biệt thứ tự. Cuối cùng, dùng thứ tự của
phần $B$ làm khóa thứ hai và trộn với phần $A$ đã được sắp xếp theo khóa thứ
nhất, ta thu được mảng hậu tố cuối cùng.

\begingroup
\small
\begin{verbatim}
#define F(x) ((x)/3+((x)%3==1?0:tb))
#define G(x) ((x)<tb?(x)*3+1:((x)-tb)*3+2)
int wa[maxn],wb[maxn],wv[maxn],ws[maxn];
int c0(int *r,int a,int b) {
    return r[a]==r[b]&&r[a+1]==r[b+1]&&r[a+2]==r[b+2];
}
int c12(int k,int *r,int a,int b) {
    if(k==2)
        return r[a]<r[b]||r[a]==r[b]&&c12(1,r,a+1,b+1);
    else
        return r[a]<r[b]||r[a]==r[b]&&wv[a+1]<wv[b+1];
}
void sort(int *r,int *a,int *b,int n,int m) {
    int i;
    for(i=0;i<n;i++) wv[i]=r[a[i]];
    for(i=0;i<m;i++) ws[i]=0;
    for(i=0;i<n;i++) ws[wv[i]]++;
    for(i=1;i<m;i++) ws[i]+=ws[i-1];
    for(i=n-1;i>=0;i--) b[--ws[wv[i]]]=a[i];
    return;
}
void dc3(int *r,int *sa,int n,int m)
{
    int i,j,*rn=r+n,*san=sa+n,ta=0,tb=(n+1)/3,tbc=0,p;
    r[n]=r[n+1]=0;
    for(i=0;i<n;i++) if(i%3!=0) wa[tbc++]=i;
    sort(r+2,wa,wb,tbc,m);
    sort(r+1,wb,wa,tbc,m);
    sort(r,wa,wb,tbc,m);
    for(p=1,rn[F(wb[0])]=0,i=1;i<tbc;i++)
    rn[F(wb[i])]=c0(r,wb[i-1],wb[i])?p-1:p++;
    if(p<tbc) dc3(rn,san,tbc,p);
    else for(i=0;i<tbc;i++) san[rn[i]]=i;
    for(i=0;i<tbc;i++) if(san[i]<tb) wb[ta++]=san[i]*3;
    if(n%3==1) wb[ta++]=n-1;
    sort(r,wb,wa,ta,m);
    for(i=0;i<tbc;i++) wv[wb[i]=G(san[i])]=i;
    for(i=0,j=0,p=0;i<ta && j<tbc;p++)
    sa[p]=c12(wb[j]%3,r,wa[i],wb[j])?wa[i++]:wb[j++];
    for(;i<ta;p++) sa[p]=wa[i++];
    for(;j<tbc;p++) sa[p]=wb[j++];
    return;
}
\end{verbatim}
\endgroup

\begin{enumerate}
  \item Trong tham số ở dòng 24, hai mảng \texttt{rn} và \texttt{san} phục vụ
  lời gọi đệ quy. \texttt{rn} lưu thông tin thứ hạng của xâu hiện tại, còn
  \texttt{san} lưu mảng hậu tố của xâu hiện tại.

  \item Dòng 25 đặt thêm $0$ vì trước khi nối hai phần có vị trí không đồng dư
  $0$ theo modulo $3$, ta cần bảo đảm độ dài mỗi phần là bội của $3$ để có thể
  sắp xếp cơ số. Với độ dài bất kỳ, thêm nhiều nhất hai số $0$ là đủ; ở đây
  làm vậy để tiện cho các bước sau.

  \item Dòng 26--29 sắp xếp cơ số các xâu con của xâu mới, mỗi xâu con gồm
  một nhóm $3$ ký tự, cuối cùng lưu mảng hậu tố vào \texttt{wb}. Chú ý rằng vị
  trí của hai mảng trung gian \texttt{wa} và \texttt{wb} được hoán đổi một lần;
  việc này làm cho kết quả sắp xếp theo chữ số cao hơn đồng thời phụ thuộc vào
  kết quả đã sắp xếp theo chữ số thấp hơn.

  \item Dòng 30--31 dùng hàm \texttt{F} để ánh xạ vị trí bắt đầu của hậu tố
  trong xâu ban đầu sang vị trí trong xâu mới. Cần nhớ rằng trong xâu mới,
  $3$ ký tự được xem như $1$ ký tự. Từ đó ta tính mảng thứ hạng
  \texttt{rn} của các ký tự trong xâu mới.

  \item Dòng 32--33: nếu trong lần sắp xếp này mọi ký tự đều không hoàn toàn
  giống nhau, nghĩa là đã phân biệt được toàn bộ thứ tự, ta trực tiếp tính
  mảng hậu tố và lưu vào \texttt{san}. Nếu vẫn tồn tại các ký tự hoàn toàn
  giống nhau, ta gọi đệ quy hàm này và kết quả vẫn nằm trong \texttt{san}. Vì
  trước mỗi lần đệ quy, độ dài xâu mới không vượt quá $2/3$ độ dài xâu ban
  đầu, đệ quy chắc chắn dừng.

  \item Dòng 34--36 dùng phần có vị trí không đồng dư $0$ theo modulo $3$ làm
  khóa thứ hai, lấy trực tiếp thông tin thứ tự từ \texttt{san} hiện tại. Cần
  chú ý trường hợp đặc biệt \texttt{n \% 3 == 1}: do không tồn tại
  \texttt{san[n]}, ta phải gán riêng. Sau đó dùng sắp xếp cơ số để trộn với
  phần đồng dư $0$ theo modulo $3$, tức khóa thứ nhất, thu được mảng hậu tố
  của phần đồng dư $0$ và lưu trong \texttt{wa}.

  \item Dòng 37 ánh xạ vị trí từng ký tự của xâu mới về xâu ban đầu bằng hàm
  \texttt{G}, từ đó thu được mảng hậu tố \texttt{wb} và mảng thứ hạng
  \texttt{wv} cho các hậu tố có vị trí không đồng dư $0$ theo modulo $3$ của
  xâu ban đầu.

  \item Dòng 38--41 trộn hai phần: phần đồng dư $0$ và phần không đồng dư $0$
  theo modulo $3$. Khi mọi phần tử của một trong hai mảng hậu tố \texttt{wa}
  hoặc \texttt{wb} đã được lấy ra, thứ tự của các phần tử còn lại trong mảng
  kia đều đứng sau, nên chỉ cần nối chúng vào cuối \texttt{sa}.
\end{enumerate}

Ví dụ trong bài là xâu \texttt{aabaaaaba}. Khi thêm ký tự canh \texttt{0}, ta
có
\[
  r=[1,1,2,1,1,1,1,2,1,0],\qquad n=10.
\]
Các vị trí không đồng dư $0$ theo modulo $3$ ban đầu là
\[
  wa=[1,2,4,5,7,8].
\]
Sau khi sắp xếp ba khóa, gán thứ hạng cho xâu mới, sắp xếp phần đồng dư $0$
và trộn hai phần, mảng hậu tố cuối cùng là
\[
  sa=[9,8,3,4,5,0,6,1,7,2].
\]

\section{Mảng height}

\subsection*{Định nghĩa ngắn gọn}

Mảng \texttt{height} là độ dài tiền tố chung dài nhất của hai hậu tố kề nhau
trong thứ tự đã sắp. Ví dụ, với hậu tố \texttt{sa[rank[i]]} của
\texttt{rank[i]} và hậu tố \texttt{sa[rank[i] - 1]} đứng ngay trước nó, tiền
tố chung dài nhất của chúng được ghi là \texttt{height[rank[i]]}, viết tắt là
$h[i]$. Nói cách khác, $h[i]$ là độ dài tiền tố chung dài nhất giữa hậu tố bắt
đầu ở vị trí $i$ trong xâu và hậu tố đứng ngay trước nó theo thứ tự từ điển.

\subsection*{Tính chất quan trọng}

Với mọi $i\ge 1$, ta có
\[
  h[i]\ge h[i-1]-1.
\]

Chứng minh: giả sử $i$ và $j$ thỏa
\[
  \operatorname{rank}[i]-\operatorname{rank}[j]=1.
\]
Khi đó độ dài tiền tố chung dài nhất của hai hậu tố này là $h[i]$. Bây giờ xét
hai hậu tố trong xâu ban đầu có vị trí bắt đầu sau $j$ và sau $i$, ký hiệu là
$j+1$ và $i+1$.

Trường hợp thứ nhất: ký tự đầu của hậu tố $i$ và hậu tố $j$ khác nhau. Khi đó
$h[i]=0$, nên hiển nhiên $h[i+1]\ge h[i]-1$.

Trường hợp thứ hai: ký tự đầu của hậu tố $i$ và hậu tố $j$ giống nhau. Khi đó
hậu tố $i+1$ và hậu tố $j+1$ lần lượt là hậu tố $i$ và hậu tố $j$ sau khi bỏ
ký tự đầu. Rõ ràng hậu tố $j+1$ đứng trước hậu tố $i+1$, và tiền tố chung dài
nhất của chúng có độ dài $h[i]-1$. Trong các hậu tố đứng trước hậu tố $i+1$
theo thứ tự sắp xếp, hậu tố kề với nó chắc chắn có tiền tố chung dài nhất với
nó; có thể hiểu là độ tương tự cao nhất. Do đó hậu tố có thứ hạng
\[
  \operatorname{rank}[i+1]-1
\]
và hậu tố có thứ hạng $\operatorname{rank}[i+1]$ có tiền tố chung dài nhất ít
nhất là $h[i]-1$. Suy ra
\[
  h[i+1]\ge h[i]-1.
\]

Lấy ví dụ xâu \texttt{aabaaaab}. Mảng \texttt{rank} của nó là
\[
  [4,6,8,1,2,3,5,7].
\]
Giả sử $i=6$ và $j=5$. Khi đó
\[
  suffix(i)=\texttt{ab},\qquad suffix(j)=\texttt{aab}.
\]
Hậu tố thứ $i+1$ là \texttt{b}, còn hậu tố thứ $j+1$ là \texttt{ab}. Hậu tố
$j+1$ đứng trước hậu tố $i+1$, đồng thời
\[
  \operatorname{rank}[i+1]=8,\qquad \operatorname{rank}[j+1]=7.
\]
Vì thứ hạng của hậu tố $j+1$ vừa đúng là thứ hạng ngay trước hậu tố $i+1$, ở
ví dụ này đạt dấu bằng:
\[
  \texttt{height[rank[7]]}\ge \texttt{height[rank[6]]}-1.
\]

\subsection*{Cài đặt}

\begingroup
\small
\begin{verbatim}
int rank[maxn],height[maxn];
void calheight(int *r,int *sa,int n)
{
    int i,j,k=0;
    for(i=1;i<=n;i++) rank[sa[i]]=i;
    for(i=0;i<n;height[rank[i++]]=k)
        for(k?k--:0,j=sa[rank[i]-1];r[i+k]==r[j+k];k++);
    return;
}
\end{verbatim}
\endgroup

\begin{enumerate}
  \item Dòng 5 dùng mảng hậu tố của xâu để tính mảng thứ hạng tương ứng.

  \item Dòng 6--7: khi \texttt{k} bằng $0$, ta bắt đầu tìm lần lượt từ ký tự
  đầu tiên. Sau đó, theo thứ tự của mảng thứ hạng, ta lần lượt tìm giá trị
  $h$ tiếp theo. Nhờ tính chất trên, lúc này không cần tìm lại từ đầu mà chỉ
  cần bắt đầu từ vị trí của ký tự thứ \texttt{k - 1}; cứ như vậy, ta thu được
  toàn bộ giá trị của mảng \texttt{height}.
\end{enumerate}

\end{document}
